\chapter{Literature Review and Theoretical Background}
\label{chap:litreview}

\section{Core Concepts and Definitions}
\label{sec:concepts}

Software architecture defines the fundamental structural organization of a system, comprising its computational components, their externally visible properties, and the relationships governing their interaction and evolution \cite{martin2017,ford2022}. Within modern web engineering, two structural styles dominate backend design: the \textit{modular monolith} and \textit{microservices}.

\begin{itemize}[leftmargin=1.5em]
  \item \textbf{Monolithic Architecture:} A software system packaged, deployed, and executed as a single unified process or tightly coupled deployment artifact \cite{richardson2018}. When structured using strict internal namespace boundaries, dependency inversion, and domain-aligned modules, it is termed a \textit{modular monolith}, achieving compile-time modularity without network serialization boundaries \cite{martin2017,newman2021}.
  \item \textbf{Microservices Architecture:} An architectural style that structures an application as a suite of small, autonomous services modeled around business capabilities, each running in its own process, owning its private data store, and communicating via lightweight mechanisms such as HTTP/REST, gRPC, or asynchronous message brokers \cite{fowler2014,dragoni2017}.
  \item \textbf{Bounded Context:} A central pattern in Domain-Driven Design (DDD) that demarcates an explicit linguistic and logical boundary within which a particular domain model is defined and applicable \cite{evans2003}. Both modular monoliths and microservices rely on bounded contexts to prevent domain model entanglement.
  \item \textbf{Change Propagation (Ripple Effect):} The extent to which modifying a requirement in one module necessitates structural edits, recompilation, or redeployment across dependent modules or services \cite{parnas1972,ford2022}.
\end{itemize}

\section{Related Work: Themes, Methods, and Evidence}
\label{sec:related}

Empirical and secondary literature comparing monolithic and microservices architectures spans three primary research strands:

\begin{itemize}[leftmargin=1.5em]
  \item \textbf{Systematic Mapping and Industry Migration Studies:} Pahl and Jamshidi \cite{pahl2016} and Soldani et al.\ \cite{soldani2018} synthesized industrial reports on microservices adoption, identifying independent deployability and fault isolation as primary gains, while documenting operational complexity, distributed testing difficulty, and infrastructure cost inflation as recurring pains. Similarly, Taibi et al.\ \cite{taibi2017} and Fritzsch et al.\ \cite{fritzsch2019} surveyed industry migrations, finding that teams frequently underestimate the DevOps maturity required to operate distributed service meshes and sagas.
  \item \textbf{Quantitative Runtime Performance Benchmarks:} Villamizar et al.\ \cite{villamizar2015} evaluated cloud deployment costs for monolithic versus microservice web applications, demonstrating that microservices can reduce infrastructure cost at extreme scale via granular auto-scaling, but incur higher baseline overhead at low-to-moderate traffic. Al-Debagy and Martinek \cite{aldebis2020}, Tapia et al.\ \cite{tapia2020}, and Blinowski et al.\ \cite{blinov2021} conducted controlled HTTP load tests showing that inter-service serialization and API gateway routing introduce measurable overhead on simple transactions, whereas decomposed read paths can outperform monolithic thread pools under concurrent load.
  \item \textbf{Socio-Technical and Evolutionary Architecture Theory:} Conway's Law \cite{conway1968} posits that system designs mirror the communication structures of the organizations that build them. Building on Parnas's information-hiding criteria \cite{parnas1972}, Fowler \cite{fowler2015monolith} and Newman \cite{newman2021} advocate a ``Monolith-First'' evolutionary strategy, warning that decomposing a system before domain boundaries stabilize results in a ``distributed monolith'' that combines network latency with tight coupling.
\end{itemize}

Despite extensive literature on large-scale cloud migrations, existing empirical studies rarely isolate the intersection of \textit{resource-constrained SME hardware} and \textit{standardized functional change propagation} across production-equivalent codebases written in the same language ecosystem.

\section{Theoretical Framework}
\label{sec:theory}

To evaluate both architectural styles rigorously, this study adopts a multi-dimensional quality model grounded in ISO/IEC~25010:2023 \cite{iso25010} and architectural trade-off analysis \cite{ford2022}. Specifically, architectural suitability $S(A)$ for an SME context is modeled across three orthogonal dimensions:

\begin{equation}
\label{eq:tradeoff}
U(A) = w_1 \cdot E_{\text{res}}(A) + w_2 \cdot P_{\text{lat}}(A) + w_3 \cdot M_{\text{evol}}(A), \quad \sum_{i=1}^{3} w_i = 1
\end{equation}

where $E_{\text{res}}(A)$ denotes runtime resource efficiency (inverse memory footprint and cold-start overhead), $P_{\text{lat}}(A)$ captures responsiveness under concurrent load (median $p_{50}$, tail $p_{95}$ latency, and error-free throughput), and $M_{\text{evol}}(A)$ represents evolutionary maintainability (localization of code edits and deployment scope during feature changes).

\begin{figure}[h]
  \centering
  \begin{tikzpicture}[
    node distance=1.2cm and 1.4cm,
    box/.style={rectangle, draw=black!80, thick, rounded corners=3pt, minimum width=3.4cm, minimum height=0.95cm, align=center, font=\small},
    hub/.style={rectangle, draw=black!90, very thick, fill=black!5, rounded corners=4pt, minimum width=3.8cm, minimum height=1.1cm, align=center, font=\small\bfseries}
  ]
    \node[hub] (iso) {ISO/IEC 25010:2023\\Architectural Quality};
    \node[box, above left=of iso] (res) {\textbf{Resource Efficiency}\\Idle \& Peak RAM, Boot Time};
    \node[box, above right=of iso] (perf) {\textbf{Performance Efficiency}\\$p_{50}$ / $p_{95}$ Latency, RPS};
    \node[box, below left=of iso] (maint) {\textbf{Maintainability}\\Files Touched, Coupling};
    \node[box, below right=of iso] (ops) {\textbf{Operability (SME)}\\Deploy Units, Infra Cost};

    \draw[->, thick] (res) -- (iso);
    \draw[->, thick] (perf) -- (iso);
    \draw[->, thick] (maint) -- (iso);
    \draw[->, thick] (ops) -- (iso);
  \end{tikzpicture}
  \caption{Theoretical evaluation framework mapping ISO/IEC 25010:2023 quality attributes to empirical metrics.}
  \label{fig:theory-framework}
\end{figure}

\section{Research Motivation and Direction}
\label{sec:motivation}

\paragraph{Purpose and Problem Significance.}
For SME engineering leads and founders, choosing microservices solely due to industry hype can exhaust limited server memory and slow down small developer teams, whereas choosing an unstructured monolith risks technical debt during rapid product iteration. Providing concrete, reproducible measurements on identical .NET reference implementations equips SME decision-makers with empirical evidence rather than anecdotal claims.

\paragraph{Identified Gap and Research Direction.}
Existing benchmarks seldom measure both runtime resource consumption and source-level change propagation on equivalent enterprise-grade .NET e-commerce platforms under constrained host resources. This thesis addresses that gap by pairing controlled Docker/\texttt{k6} runtime benchmarks with five standardized code-evolution scenarios across \texttt{eShopOnWeb} and \texttt{eShopOnContainers}, corroborated by practitioner interviews.

\paragraph{Success Criteria.}
The research direction is validated if: (1) both reference architectures execute reproducibly under identical traffic profiles with $0.00\%$ HTTP error rates; (2) hypotheses H1--H3 are quantitatively confirmed or refuted with exact telemetry logs; and (3) the resulting decision matrix provides unambiguous threshold rules for SME teams.

\section{Chapter Summary}
\label{sec:ch1-summary}
This chapter established the formal definitions of modular monolithic and microservices architectures, synthesized empirical findings across performance and migration literature, formulated a three-dimensional evaluation lens grounded in ISO/IEC~25010:2023, and defined the empirical gap and success criteria guiding Chapters~2 through~5.
